%% ===========================================================================
%%  preamble.tex — shared design system
%%
%%  Loaded by BOTH acl_latex.tex (the paper) and design_kit.tex (the pattern
%%  gallery), so anything defined here is guaranteed to look the same in both.
%%  Never put content here; content lives in sections/.
%%
%%  Hard constraints this file respects (see formatting.md, paper/guideline.md):
%%   * acl.sty is official and UNMODIFIED — all of our styling happens here.
%%   * acl.sty already loads: geometry, caption (font=10pt), natbib, hyperref,
%%     and (in review mode) lineno. Loading those again with options would
%%     raise an option clash, so we only ever \captionsetup{...} them.
%%   * Body font must stay Times at 11pt with the style's leading; we adjust
%%     float/caption spacing only, never \baselineskip or the type area.
%% ===========================================================================


%% ---------------------------------------------------------------------------
%%  0. Standard includes (verbatim from the official acl_latex.tex)
%% ---------------------------------------------------------------------------

\usepackage{times}
\usepackage{latexsym}
\usepackage[T1]{fontenc}     % correct rendering/hyphenation of accented words
\usepackage[utf8]{inputenc}
\usepackage{microtype}       % better line breaking; saves a surprising amount
\usepackage{inconsolata}     % typewriter face for code and prompts


%% ---------------------------------------------------------------------------
%%  1. Packages
%% ---------------------------------------------------------------------------

% -- tables ----------------------------------------------------------------
\usepackage{booktabs}       % \toprule \midrule \bottomrule — the only rules we use
\usepackage{multirow}       % vertically merged cells
\usepackage{makecell}       % line breaks inside a cell: \makecell{a\\b}
\usepackage{tabularx}       % X columns that soak up the remaining width
\usepackage{threeparttable} % table notes that inherit the table's width
\usepackage{supertabular}   % the ONLY long-table package that breaks across
                            % columns in two-column mode (longtable cannot)
\usepackage{siunitx}        % S columns for decimal-aligned numbers.
                            % Only v2+v3-portable syntax is used: S[table-format=...]
                            % and \num{...}. Do not add \sisetup keys that were
                            % renamed in siunitx v3 (e.g. table-number-alignment).
\usepackage{adjustbox}      % max width=\linewidth — shrink-to-fit, see \fitbox

% -- figures ---------------------------------------------------------------
\usepackage{graphicx}
\usepackage{subcaption}     % subfigures; must come after caption (acl.sty loaded it)

% -- algorithms ------------------------------------------------------------
\usepackage{algorithm}      % the float
\usepackage{algpseudocode}  % the pseudocode body

% -- code and boxes --------------------------------------------------------
\usepackage{listings}
\usepackage{xcolor}
\usepackage{tcolorbox}
\tcbuselibrary{breakable,skins,listings}

% -- symbols, math, lists --------------------------------------------------
\usepackage{amsmath}
\usepackage{amssymb}
\usepackage{pifont}         % \ding{51} checkmark for the vote profiles
\usepackage{enumitem}

% -- cross references (LAST: cleveref must see hyperref and all float types) --
\usepackage[capitalise,nameinlink]{cleveref}


%% ---------------------------------------------------------------------------
%%  2. Spacing
%%
%%  ACL papers are page-budget-bound (8 pages of content). These are the only
%%  legitimate savings: float separation and caption skips. Everything here is
%%  deliberately mild — if a draft overruns, cut text, do not shrink type.
%% ---------------------------------------------------------------------------

\setlength{\textfloatsep}{10pt plus 2pt minus 2pt}   % float <-> text, top/bottom
\setlength{\floatsep}{10pt plus 2pt minus 2pt}       % float <-> float
\setlength{\intextsep}{10pt plus 2pt minus 2pt}      % around [h] floats
\setlength{\dblfloatsep}{10pt plus 2pt minus 2pt}    % same, spanning floats
\setlength{\dbltextfloatsep}{12pt plus 2pt minus 2pt}

\captionsetup{skip=5pt}                              % caption <-> content
\captionsetup[sub]{font=small,skip=3pt}

% Table typography: booktabs wants no vertical rules and a little air.
\setlength{\tabcolsep}{4.5pt}
\renewcommand{\arraystretch}{1.06}

% Compact lists (ACL drafts overflow on list padding more than anywhere else).
\setlist[itemize]{leftmargin=1.2em,topsep=2pt,itemsep=1pt,parsep=0pt}
\setlist[enumerate]{leftmargin=1.4em,topsep=2pt,itemsep=1pt,parsep=0pt}
\setlist[description]{leftmargin=0pt,topsep=2pt,itemsep=1pt,parsep=0pt,style=nextline}

% Float placement discipline: ACL floats belong at the top (or bottom) of a
% column/page, never mid-text. Defaults nudge LaTeX in that direction; write
% [t] explicitly on every float anyway.
\renewcommand{\topfraction}{0.9}
\renewcommand{\bottomfraction}{0.7}
\renewcommand{\textfraction}{0.1}
\renewcommand{\floatpagefraction}{0.75}
\setcounter{topnumber}{3}
\setcounter{dbltopnumber}{3}


%% ---------------------------------------------------------------------------
%%  3. Colours
%%
%%  One restrained palette, reused by boxes, code, and figures. Every colour
%%  must survive greyscale printing and be distinguishable to colour-blind
%%  readers, so hue never carries meaning alone (boxes also differ by title).
%% ---------------------------------------------------------------------------

\definecolor{cpInk}{HTML}{1A1A1A}       % text on light fills
\definecolor{cpRule}{HTML}{9AA0A6}      % box frames, hairlines
\definecolor{cpWash}{HTML}{F4F5F7}      % neutral fill (system prompts, code)
\definecolor{cpBlue}{HTML}{2B5C8A}      % user/input, primary accent
\definecolor{cpGreen}{HTML}{2E6B4F}     % model output
\definecolor{cpAmber}{HTML}{8A6A1F}     % judge / rubric
\definecolor{cpPlum}{HTML}{6B3F7A}      % examples, retrieved evidence
\definecolor{cpRed}{HTML}{A33A32}       % failures, drafting notes


%% ---------------------------------------------------------------------------
%%  4. Tables — the three patterns
%%
%%  (a) column-width table   : table  + \fitbox   -> shrinks only if needed
%%  (b) page-width table     : table* + \fitbox   -> spans both columns
%%  (c) multi-page table     : \begin{cptablelong} (supertabular), appendix only
%%
%%  Rules for all three: caption ABOVE the tabular, booktabs rules only,
%%  no vertical lines, units in the header, \best to mark winners.
%% ---------------------------------------------------------------------------

% Shrink-to-fit wrapper. \linewidth resolves to the column inside `table` and
% to the full text width inside `table*`, so the SAME wrapper serves both and
% a table can never overflow its container regardless of content. It only ever
% scales down (max width), so small tables keep the body font size.
\newenvironment{fitbox}
  {\begin{adjustbox}{max width=\linewidth,center}}
  {\end{adjustbox}}

% Font size for table bodies. Change here, not per table.
\newcommand{\tabfont}{\small}

% Result marking (bold = best, underline = runner-up; both readable in grey).
\newcommand{\best}[1]{\textbf{#1}}
\newcommand{\second}[1]{\underline{#1}}
\newcommand{\na}{\textendash}                    % measured, not applicable
\newcommand{\nr}{$\emptyset$}                    % not run (matches vote profile)

% Significance markers for comparison tables.
\newcommand{\sigone}{\textsuperscript{*}}
\newcommand{\sigtwo}{\textsuperscript{**}}
\newcommand{\sigthree}{\textsuperscript{***}}

% Header cell that wraps in a narrow column: \hcell{false alarm\\rate}
\newcommand{\hcell}[1]{\makecell[bc]{#1}}

% Group heading spanning several columns, with a rule underneath.
\newcommand{\gcol}[2]{\multicolumn{#1}{c}{#2}}

% Multi-page table in a single column (appendix). supertabular, not longtable:
% longtable is illegal in two-column mode. Caption goes in \topcaption before.
\newenvironment{cptablelong}[1]
  {\tabfont\begin{supertabular}{#1}}
  {\end{supertabular}}


%% ---------------------------------------------------------------------------
%%  5. Figures
%%
%%  Export conventions (keep figures honest and legible at print size):
%%   * vector PDF, never PNG, unless the plot has raster content
%%   * single column: 3.15 in wide  ·  spanning: 6.5 in wide
%%   * label/tick font 8 pt so it matches this caption size after \includegraphics
%%     at width=\columnwidth (i.e. do NOT scale figures in LaTeX)
%%   * no in-figure titles — the caption is the title
%% ---------------------------------------------------------------------------

% Grey placeholder that reserves the real figure's footprint. Lets us lay out
% and page-count before any plot exists: \figph{\columnwidth}{4.2cm}{caption id}
\newcommand{\figph}[3]{%
  \setlength{\fboxsep}{0pt}%
  \fcolorbox{cpRule}{cpWash}{%
    \parbox[c][#2][c]{#1}{\centering\small\color{cpRule}\ttfamily #3}}}


%% ---------------------------------------------------------------------------
%%  6. Algorithms
%%
%%  (a) single column : \begin{algorithm}[t]
%%  (b) page width    : \begin{algorithm*}[t]  (float, cannot break)
%%  (c) unbounded     : \begin{cpalgorithm}    (not a float: breaks across
%%                       columns and pages, for long appendix pseudocode)
%% ---------------------------------------------------------------------------

\floatname{algorithm}{Algorithm}
\algrenewcommand\algorithmicrequire{\textbf{Input:}}
\algrenewcommand\algorithmicensure{\textbf{Output:}}
\algrenewcommand{\algorithmiccomment}[1]{\hfill\small\textcolor{cpRule}{\# #1}}
\algnewcommand{\LineComment}[1]{\State\small\textcolor{cpRule}{\# #1}}
\newcommand{\algfont}{\small}

% Non-floating algorithm: identical look, but the body may split. Numbering
% shares the `algorithm` counter, so references stay consistent with (a)/(b).
\newenvironment{cpalgorithm}[1][]
  {\par\addvspace{8pt}\refstepcounter{algorithm}%
   \noindent\rule{\linewidth}{0.8pt}\par\vspace{2pt}%
   \noindent\textbf{\algorithmname~\thealgorithm}%
   \ifx\relax#1\relax\else\quad #1\fi\par\vspace{1pt}%
   \noindent\rule{\linewidth}{0.4pt}\par\vspace{2pt}\algfont}
  {\par\vspace{2pt}\noindent\rule{\linewidth}{0.8pt}\par\addvspace{8pt}}
% \algorithmname is defined by \floatname above; alias for the environment:
\providecommand{\algorithmname}{Algorithm}


%% ---------------------------------------------------------------------------
%%  7. Code
%%
%%  Two shapes only:
%%   \begin{pycode}[caption]  — verbatim Python, breaks across columns/pages
%%   \begin{pycodenum}[...]   — same with line numbers (for text that refers
%%                              to specific lines)
%%  Page-width code must sit inside a figure* float (floats cannot break).
%%  Inline: \code{EnsembleDetector} for identifiers in running text.
%% ---------------------------------------------------------------------------

\newcommand{\code}[1]{\texttt{\small #1}}

\lstdefinestyle{cppython}{
  language=Python,
  basicstyle=\scriptsize\ttfamily\color{cpInk},
  keywordstyle=\color{cpBlue}\bfseries,
  commentstyle=\color{cpGreen},
  stringstyle=\color{cpPlum},
  identifierstyle=\color{cpInk},
  showstringspaces=false,
  breaklines=true,
  breakatwhitespace=false,
  postbreak=\mbox{\textcolor{cpRule}{$\hookrightarrow$}\space},
  tabsize=4,
  columns=fullflexible,
  keepspaces=true,
  upquote=true,
  literate={-}{-}1 {'}{'}1,       % keep hyphens/quotes copy-pasteable
}
\lstdefinestyle{cpjson}{
  basicstyle=\scriptsize\ttfamily\color{cpInk},
  stringstyle=\color{cpPlum},
  showstringspaces=false,
  breaklines=true,
  columns=fullflexible,
  keepspaces=true,
  upquote=true,
}
\lstdefinestyle{cpshell}{
  language=bash,
  basicstyle=\scriptsize\ttfamily\color{cpInk},
  keywordstyle=\color{cpBlue},
  commentstyle=\color{cpGreen},
  showstringspaces=false,
  breaklines=true,
  columns=fullflexible,
}

% Shared geometry for every box in the document (boxes read as one family).
\tcbset{cpbox/.style={
  enhanced, breakable,
  boxrule=0.5pt, arc=1.5pt,
  left=5pt, right=5pt, top=4pt, bottom=4pt,
  fonttitle=\small\bfseries, coltitle=cpInk,
  attach boxed title to top left={yshift=-2mm,xshift=4mm},
  boxed title style={boxrule=0.3pt, arc=1pt},
}}

% Grey frame + wash fill, shared by all code boxes.
\tcbset{cpcode/.style={cpbox, listing only,
  colback=cpWash, colframe=cpRule,
  boxed title style={colback=cpWash!60!white, colframe=cpRule}}}

\newtcblisting{pycode}[1][Python]{cpcode, listing style=cppython, title={#1}}

\newtcblisting{pycodenum}[1][Python]{cpcode, title={#1},
  listing options={style=cppython, numbers=left,
                   numberstyle=\tiny\color{cpRule}, numbersep=6pt}}

% Untitled variant for two-line snippets where a title box is overkill.
\newtcblisting{pycodeplain}{cpcode, listing style=cppython}

\newtcblisting{shellcode}[1][Shell]{cpcode, listing style=cpshell, title={#1}}

\newtcblisting{jsoncode}[1][JSON]{cpcode, listing style=cpjson, title={#1}}


%% ---------------------------------------------------------------------------
%%  8. LLM prompt boxes
%%
%%  Prompts are data, not prose: they contain braces, underscores, and JSON, so
%%  the default is a VERBATIM box (nothing to escape, copy-pasteable, breaks
%%  across columns and pages):
%%     \begin{sysprompt} ... \end{sysprompt}     system prompt
%%     \begin{userprompt} ... \end{userprompt}   user / augmented prompt
%%     \begin{modelout} ... \end{modelout}       model response
%%     \begin{judgeprompt} ... \end{judgeprompt} LLM-as-judge instructions
%%     \begin{evidence} ... \end{evidence}       retrieved chunk / document
%%
%%  Use the -fmt variants only when a prompt needs typeset emphasis (e.g.
%%  \ph{placeholder} highlighting in the paper's method section); there,
%%  every special character must be escaped by hand.
%% ---------------------------------------------------------------------------

% Placeholder inside a formatted prompt: \ph{break\_date}
\newcommand{\ph}[1]{{\normalfont\itshape\color{cpBlue}\{\{#1\}\}}}

\tcbset{cpprompt/.style={cpbox,
  fontupper=\scriptsize\ttfamily,
  before upper={\parindent0pt\parskip3pt\raggedright\sloppy}}}

% Verbatim prompt text: no language keywords, wrap long lines, mark the wrap.
\lstdefinestyle{cpprompttext}{
  basicstyle=\scriptsize\ttfamily\color{cpInk},
  breaklines=true, columns=fullflexible, keepspaces=true, upquote=true,
  showstringspaces=false,
  postbreak=\mbox{\textcolor{cpRule}{$\hookrightarrow$}\space},
}
\tcbset{cppromptverb/.style={cpprompt, listing only, listing style=cpprompttext}}

\newtcblisting{sysprompt}[1][System prompt]{cppromptverb,
  colback=cpWash, colframe=cpRule,
  boxed title style={colback=cpWash!60!white, colframe=cpRule}, title={#1}}

\newtcblisting{userprompt}[1][User prompt]{cppromptverb,
  colback=cpBlue!4!white, colframe=cpBlue!55!white,
  boxed title style={colback=cpBlue!12!white, colframe=cpBlue!55!white}, title={#1}}

\newtcblisting{modelout}[1][Model output]{cppromptverb,
  colback=cpGreen!4!white, colframe=cpGreen!55!white,
  boxed title style={colback=cpGreen!12!white, colframe=cpGreen!55!white}, title={#1}}

\newtcblisting{judgeprompt}[1][Judge prompt]{cppromptverb,
  colback=cpAmber!5!white, colframe=cpAmber!55!white,
  boxed title style={colback=cpAmber!14!white, colframe=cpAmber!55!white}, title={#1}}

% Formatted (non-verbatim) counterparts — support \ph{} and \textbf{}.
\newtcolorbox{syspromptfmt}[1][System prompt]{cpprompt,
  colback=cpWash, colframe=cpRule,
  boxed title style={colback=cpWash!60!white, colframe=cpRule}, title={#1}}
\newtcolorbox{userpromptfmt}[1][User prompt]{cpprompt,
  colback=cpBlue!4!white, colframe=cpBlue!55!white,
  boxed title style={colback=cpBlue!12!white, colframe=cpBlue!55!white}, title={#1}}
\newtcolorbox{modeloutfmt}[1][Model output]{cpprompt,
  colback=cpGreen!4!white, colframe=cpGreen!55!white,
  boxed title style={colback=cpGreen!12!white, colframe=cpGreen!55!white}, title={#1}}

% Retrieved document / corpus excerpt (RAG sections). Prose, so not verbatim.
\newtcolorbox{evidence}[1][Retrieved document]{cpbox,
  fontupper=\small, before upper={\parindent0pt\parskip3pt},
  colback=cpPlum!4!white, colframe=cpPlum!55!white,
  boxed title style={colback=cpPlum!12!white, colframe=cpPlum!55!white}, title={#1}}

% Qualitative example (input -> output pair discussed in the text).
\newtcolorbox{examplebox}[1][Example]{cpbox,
  fontupper=\small, before upper={\parindent0pt\parskip3pt},
  colback=white, colframe=cpInk!35!white,
  boxed title style={colback=cpInk!8!white, colframe=cpInk!35!white}, title={#1}}


%% ---------------------------------------------------------------------------
%%  9. Domain vocabulary
%%
%%  Every recurring name gets a macro, so a rename is one edit and the paper
%%  can never disagree with itself about a member name or a symbol.
%% ---------------------------------------------------------------------------

% Tool name. REVIEW MODE: the submission is anonymous and the repository link
% must be an anonymized one — check that naming the package does not identify
% the authors (a public PyPI/GitHub project with the same name would).
\newcommand{\pkg}{\textsc{whychange}}

% Ensemble members (exactly the identifiers used in the code and the tables).
\newcommand{\mMosum}{\code{mosum}}
\newcommand{\mPeltL}{\code{pelt\_l2}}
\newcommand{\mPeltN}{\code{pelt\_normal}}
\newcommand{\mKernel}{\code{kernel\_rbf}}
\newcommand{\mBaiPerron}{\code{bai\_perron}}
\newcommand{\mZA}{\code{zivot\_andrews}}
\newcommand{\mProphet}{\code{prophet\_trend}}

% Three-state vote profile: fired / ran-but-silent / not-run.
\newcommand{\vfired}{\textcolor{cpGreen}{\ding{51}}}
\newcommand{\vsilent}{\textcolor{cpRule}{\textendash}}
\newcommand{\vnotrun}{$\emptyset$}
\newcommand{\voteLegend}{%
  \vfired~fired \quad \vsilent~ran, silent \quad \vnotrun~not applicable}

% Notation. Keep these the single source of truth for the maths.
\newcommand{\nobs}{n}                                  % series length
\newcommand{\ts}{\ensuremath{y_{1:\nobs}}}             % the series
\newcommand{\pmem}[1]{\ensuremath{p_{#1}}}             % per-member p-value
\newcommand{\pcomb}{\ensuremath{p_{\mathrm{comb}}}}    % Cauchy-combined p
\newcommand{\epsclust}{\ensuremath{\varepsilon}}       % clustering tolerance
\newcommand{\alphalev}{\ensuremath{\alpha}}            % acceptance level
\newcommand{\blocklen}{\ensuremath{\ell}}              % permutation block length
\newcommand{\nperm}{\ensuremath{B}}                    % number of permutations
\DeclareMathOperator*{\argmax}{arg\,max}
\DeclareMathOperator*{\argmin}{arg\,min}

% Prose shorthands.
\newcommand{\ie}{i.e.,\ }
\newcommand{\eg}{e.g.,\ }
\newcommand{\etal}{et al.\ }


%% ---------------------------------------------------------------------------
%%  10. Drafting aids — MUST be silenced before submission
%%
%%  \draftmodefalse turns every marker below into a no-op, so a stray \todo
%%  cannot reach a reviewer. Flip it in acl_latex.tex.
%% ---------------------------------------------------------------------------

\newif\ifdraftmode
\draftmodetrue            % set \draftmodefalse in acl_latex.tex before submitting

\newcommand{\todo}[1]{\ifdraftmode\textcolor{cpRed}{\textbf{[TODO: #1]}}\fi}
\newcommand{\note}[1]{\ifdraftmode\textcolor{cpBlue}{\textbf{[NOTE: #1]}}\fi}
% A number that still has to be regenerated from the final evaluation run:
% \prov{0.84} is underlined while drafting and plain in the submission build,
% so no placeholder figure can silently survive into the PDF.
\newcommand{\prov}[1]{\ifdraftmode\textcolor{cpRed}{\underline{#1}}\else#1\fi}
% Citation whose bibliographic details are not yet verified against the source
% (EACL screens fabricated references — nothing ships with \unverified live).
\newcommand{\unverified}[1]{\ifdraftmode\textcolor{cpRed}{[verify: #1]}\fi}


% Build-mode test usable in the document body (acl.sty's own flag is internal):
% \ifanonymous ... \else ... \fi  -> true in the [review] build.
\makeatletter
\newif\ifanonymous
\ifacl@anonymize\anonymoustrue\else\anonymousfalse\fi
\makeatother


%% ---------------------------------------------------------------------------
%%  11. Cross-reference names
%% ---------------------------------------------------------------------------

\crefname{algorithm}{Algorithm}{Algorithms}
\Crefname{algorithm}{Algorithm}{Algorithms}
\crefname{section}{Section}{Sections}
\Crefname{section}{Section}{Sections}
\crefname{appsec}{Appendix}{Appendices}
\crefformat{footnote}{#2\footnotemark[#1]#3}
